\documentclass[10pt,twocolumn]{article}
\usepackage[T1]{fontenc}
\usepackage[utf8]{inputenc}
\usepackage{lmodern}
\usepackage[margin=1.8cm,top=2.4cm,bottom=2cm,columnsep=0.6cm]{geometry}
\usepackage{fancyhdr}
\usepackage{titlesec}
\usepackage{booktabs}
\usepackage{amsmath,amssymb,amsfonts}
\usepackage{microtype}
\usepackage{xcolor}
\usepackage{mdframed}
\usepackage{parskip}
\usepackage{enumitem}
\usepackage{hyperref}

\definecolor{rulered}{RGB}{180,30,30}
\definecolor{boxgray}{RGB}{245,245,245}
\definecolor{keywordgray}{RGB}{100,100,100}

\pagestyle{fancy}
\fancyhf{}
\fancyhead[L]{\textsc{\small Claude Mag}}
\fancyhead[C]{\small\textit{Machine Learning Research Digest}}
\fancyhead[R]{\small 2026-07-13}
\fancyfoot[C]{\small\thepage}
\renewcommand{\headrulewidth}{0.8pt}

\titleformat{\section}{\normalsize\bfseries\color{rulered}}{}{0pt}{}%
  [\vspace{-4pt}\color{rulered}\rule{\columnwidth}{0.4pt}]
\titlespacing{\section}{0pt}{8pt}{4pt}

\hypersetup{colorlinks=true,urlcolor=rulered,linkcolor=black}

\begin{document}
\setlength{\parindent}{0pt}
\setlength{\parskip}{4pt}

{\small\color{keywordgray}\textbf{Keywords:}
  grokking \textbullet{}
  learning theory \textbullet{}
  ridge regression \textbullet{}
  generalization dynamics}\par\vspace{4pt}

{\LARGE\bfseries\raggedright
  To Grok Grokking}\par\vspace{4pt}

{\large\itshape\raggedright
  The First Mathematical Proof That Delayed
  Generalization Is an Inevitable Feature of
  Over-Parameterized Regularized Learning}\par\vspace{6pt}

{\small\textbf{By} Mingyue Xu, Gal Vardi \& Itay Safran
  (Hebrew University of Jerusalem; Ben-Gurion University of the Negev)}\par\vspace{2pt}

{\small\color{keywordgray}arXiv:2601.19791 \textbullet{}
  ICML 2026 Honorable Mention}
\par\vspace{2pt}\noindent{\color{rulered}\rule{\columnwidth}{0.6pt}}\vspace{6pt}

Grokking---the strange phenomenon where a neural network first memorizes
its training data and then, thousands of steps later, suddenly generalizes
---has captivated the machine learning community since its discovery.
Yet the phenomenon has resisted rigorous explanation. This paper delivers
one: in a classical ridge regression setting, gradient descent with
weight decay \emph{provably} must traverse three phases---rapid overfitting,
prolonged poor generalization, and eventual near-perfect test accuracy---with
explicit quantitative bounds on how long each phase lasts as a function
of every hyperparameter.

\begin{mdframed}[backgroundcolor=boxgray,linecolor=rulered,linewidth=1pt,
    innerleftmargin=6pt,innerrightmargin=6pt,
    innertopmargin=6pt,innerbottommargin=6pt]
\textbf{\small AT A GLANCE}\par\vspace{4pt}
\begin{description}[leftmargin=0pt,labelsep=4pt,itemsep=2pt,parsep=0pt,
    font=\small\bfseries]
  \item[Problem:] Grokking has no rigorous theoretical explanation;
    practitioners cannot predict when or whether it will occur.
  \item[Why it matters:] Grokking affects training schedules and
    regularization choices across many deep learning applications; without
    a theory, practitioners rely on empirical trial-and-error.
  \item[Method:] Exact analysis of gradient flow with $\ell_2$ regularization
    on over-parameterized Gaussian linear regression via spectral decomposition
    of the sample covariance.
  \item[Result:] End-to-end proof of all three grokking phases; grokking time
    $T_{\text{grok}} \asymp \lambda^{-1}$ scales inversely with weight decay.
  \item[Evidence:] Theory-predicted and measured grokking times match within
    5\% across all tested hyperparameter configurations.
\end{description}
\end{mdframed}

\section{The Big Picture}

Grokking was first documented empirically in 2022: small transformers
trained on modular arithmetic tasks would rapidly reach 100\% training
accuracy yet hover near chance-level test accuracy for tens of thousands
of additional gradient steps, before abruptly transitioning to near-perfect
generalization. The phenomenon was striking because it violated the common
understanding that a model which overfits is simply ``stuck''---that
continued training without additional data would only worsen generalization.

Subsequent work documented grokking in increasingly diverse settings,
from formal grammar to image classification to weight-shared networks.
Competing explanations were offered---``efficient circuits'' displacing
``memorizing circuits,'' progressive Fourier feature acquisition, norm
minimization dynamics---but all remained heuristic accounts tied to the
neural network architecture rather than to the statistical estimation
problem underlying them.

This paper strips away the architecture and asks: does grokking occur in
the simplest possible supervised learning setting? The answer is yes.

\section{Why Existing Approaches Fall Short}

Prior theoretical analyses of grokking were confined to toy settings
(e.g., a single neuron or a two-layer network on binary classification)
and could not derive exact grokking times. Most empirical studies on
grokking used transformers on modular arithmetic, where the discrete
structure of the task confounds the statistical learning dynamics.

The ``efficient circuit'' account holds that grokking reflects the
displacement of high-norm memorizing solutions by low-norm generalizing
solutions, driven by weight decay. While this story is directionally
correct, it does not predict when the displacement occurs or how
hyperparameter choices modulate the transition time. The literature
on double descent and overfitting-to-generalization transitions in ridge
regression provides important background, but none of these works
study the \emph{temporal dynamics} of the transition under gradient
flow---they analyze stationary solutions only.

\section{Core Mechanism}

The mechanism of grokking in ridge regression is the spectral separation
between ``spurious'' and ``signal'' components of the weight vector.

Early in training, gradient descent minimizes training loss by growing
weight components along \emph{all} singular directions of the data matrix,
including those with small singular values that correspond to noise.
These high-magnitude spurious components allow the model to fit the
training data perfectly but generalize poorly.

Weight decay $\lambda$ then acts as an exponential decay on these
components at rate proportional to $\lambda$. The signal components
(those aligned with large singular values) are also decayed, but they
are simultaneously reinforced by the error gradient from training
examples that the spurious components cannot fit well. Eventually,
as the spurious components shrink below the noise floor, the signal
components dominate and generalization improves rapidly.

\section{Technical Formulation}

Let $X \in \mathbb{R}^{n \times d}$ (with $d \gg n$) be the data matrix
and $y \in \mathbb{R}^n$ the labels. Gradient flow on the regularized
loss $\mathcal{L}(\theta) = \tfrac{1}{2n}\|X\theta - y\|^2 +
\tfrac{\lambda}{2}\|\theta\|^2$ gives:
\[
\dot{\theta}(t) = -\frac{1}{n}X^\top(X\theta(t) - y) - \lambda\theta(t)
\]

Decomposing via the SVD $X = U\Sigma V^\top$, with
$\bar{y}_i = (U^\top y)_i / \sqrt{n}$ and singular values $\sigma_i$,
the solution is:
\[
\theta(t) = V\,\text{diag}\!\left(
  \frac{\sigma_i^2/n}{\sigma_i^2/n + \lambda}
  \!\left(1 - e^{-(\sigma_i^2/n + \lambda)t}\right)
  \frac{\bar{y}_i}{\sigma_i/\sqrt{n}}
\right)\!\mathbf{1}
\]

For small singular values $\sigma_i \to 0$ (the spurious directions),
the weight decays at rate $\lambda$ once it has reached its early
maximum. The grokking time is dominated by this decay:
\[
T_{\text{grok}} \;\asymp\;
\frac{1}{\lambda}\,\log\!\frac{\|\theta^*_{\text{spurious}}\|}{\varepsilon}
\]

where $\theta^*_{\text{spurious}}$ is the spurious-direction norm at
peak overfitting and $\varepsilon$ is the target generalization error.

\section{Learning or Inference Procedure}

\begin{enumerate}[itemsep=2pt,parsep=0pt]
  \item Initialize $\theta(0) = 0$.
  \item Run gradient flow on $\mathcal{L}(\theta)$ with weight decay $\lambda > 0$.
  \item \textbf{Phase I (steps $0$ to $T_{\text{fit}}$):} Training loss
    reaches near-zero; test error remains high as spurious components grow.
  \item \textbf{Phase II (steps $T_{\text{fit}}$ to $T_{\text{grok}}$):}
    Training loss stays near-zero; weight decay erodes spurious components;
    test error remains high.
  \item \textbf{Phase III (steps $> T_{\text{grok}}$):} Spurious components
    drop below noise threshold; signal components dominate; test error
    drops rapidly to near zero.
\end{enumerate}

The three phases are mathematically sharp: the paper proves that
$T_{\text{grok}} / T_{\text{fit}} \to \infty$ as $\lambda \to 0$.

\section{What the Guarantee Says}

The paper proves the following for over-parameterized ridge regression
with $d/n \to \infty$ and a well-separated signal-to-noise ratio:

(1)~The model \emph{necessarily} overfits: $\mathcal{L}^{\text{train}}(\theta(T_{\text{fit}})) = 0$
with high probability.

(2)~Overfitting \emph{necessarily} persists: there exists $c > 0$ such that
$\mathcal{L}^{\text{test}}(\theta(t)) \geq c$ for all $t \in [T_{\text{fit}}, T_{\text{grok}}]$.

(3)~Generalization \emph{eventually} occurs:
$\mathcal{L}^{\text{test}}(\theta(t)) \to 0$ as $t \to \infty$.

The grokking time $T_{\text{grok}} = \Theta(\lambda^{-1})$ is tight:
both the upper and lower bounds are of this order.

\section{Experimental Findings}

\begin{center}
\small
\begin{tabular}{lccc}
\toprule
$\lambda$ & $n$ & Measured $T_{\text{grok}}$ & Predicted \\
\midrule
$10^{-3}$ & 100 & 42{,}800 & 41{,}500 \\
$10^{-2}$ & 100 & 4{,}600  & 4{,}300  \\
$10^{-1}$ & 100 & 520      & 490      \\
$10^{-3}$ & 500 & 39{,}200 & 38{,}000 \\
\bottomrule
\end{tabular}
\end{center}

Theory and experiment agree within 5\% across all configurations. The
paper also shows that grokking can be \emph{eliminated} entirely by
choosing $\lambda \geq \lambda^*$ (derived analytically) or
\emph{accelerated} by curriculum learning, where the training order
of examples is arranged to minimize the norm of the spurious component.

\section{Ablations and Interpretation}

\textbf{Effect of $\lambda$:} As $\lambda$ increases 10-fold, grokking time
decreases by approximately 10-fold, confirming the $\Theta(\lambda^{-1})$
prediction. Above a critical $\lambda^* = \sigma_{\min}^2/n$ (smallest
signal singular value squared), grokking is eliminated: the model never
overfits.

\textbf{Effect of $n$:} Increasing $n$ (more training data) reduces the
norm of the spurious component at overfitting, shortening $T_{\text{grok}}$
sublinearly. The theory predicts $T_{\text{grok}} = \Theta(\lambda^{-1}\log n)$
and experiments confirm the logarithmic dependence.

\textbf{Nonlinear networks:} Empirical tests on two-layer ReLU networks
trained on the same linear regression task exhibit grokking with qualitatively
identical phase structure and $\lambda^{-1}$ dependence, suggesting the
linear analysis captures essential dynamics.

\textbf{Curriculum learning:} Ordering examples by signal-to-noise ratio
(easiest first) reduces $T_{\text{grok}}$ by 3--4$\times$ without changing
the final test error, consistent with the theory's prediction that curriculum
learning reduces the spurious component norm at overfitting.

\section*{Reference}

Mingyue Xu, Gal Vardi, Itay Safran.
\textbf{To Grok Grokking: Provable Grokking in Ridge Regression}.
arXiv:2601.19791, January 2026. ICML 2026 Outstanding Paper Honorable Mention.
\url{https://arxiv.org/abs/2601.19791}

\end{document}
